\section{Results}

\subsection{\Q1: How many attacks does the manifest block?}
The hand-written manifest denies at least one tool the injection needs in \N{manBlk}\% of the \N{nPairs} pairs (95\%
interval \N{manBlkCI}\%; Table~\ref{tab:oracle}, Fig.~\ref{fig:teaser}). The rate depends strongly on the suite:
\N{manBlkworkspace}\% in workspace, \N{manBlkslack}\% in Slack, \N{manBlktravel}\% in travel and only
\N{manBlkbanking}\% in banking. The result does not hinge on annotation D1. On the \N{nPairsVone} pairs of the original
v1 tasks it is \N{manBlkVone}\% (ceiling \N{otoolBlkVone}\%). Ignoring the clean-up \tool{delete\_email} of the
annotated tasks gives \N{manBlkStrict}\% (ceiling \N{otoolBlkStrict}\%).

The verb lists matter. With keywords taken mechanically from the tool name and docstring, the manifest blocks
\N{autoBlk}\% (\N{autoBlkCI}\%) and keeps fewer benign tasks (\N{autoPres}\% against \N{manPres}\%). Docstring words
such as ``money'', ``transaction'' or ``email'' name objects, so they match read-only requests and open sensitive
categories for no reason. In travel the auto manifest blocks \N{autoBlktravel}\%, against \N{manBlktravel}\% for the
hand-written one.

\subsection{\Q2: How many benign tasks does it keep?}
\N{manPres}\% of the \N{nUsers} user tasks keep every ground-truth call allowed (\N{manPresCI}\%). Only \N{manDenied} of
\N{nBenignCalls} benign calls are denied, but they fall in \N{nLost} tasks: \N{nLostslack} in Slack, \N{nLostworkspace}
in workspace and \N{nLostbanking} in banking. They are of three kinds. Five tasks \emph{delegate} the instruction to
content: ``do the actions specified in the email from David'', ``read the file \texttt{landlord-notices.txt} and follow
the instructions'', ``do all the tasks on my TODO list at \texttt{www.company-todo-list.com}''. The tools such a task
needs are named only in the content it points to, and that content arrives through exactly the channel an injection
uses. A manifest built from the user's words cannot serve these tasks without trusting that channel. Three Slack tasks
need to open a web page that the instruction mentions only indirectly (``summarise the article that Bob posted''), and
one says ``message them'', a verb we left out because ``message'' is far more often a noun. Travel, the suite whose
instructions name their actions most plainly (``book'', ``reserve''), keeps all of its tasks.

\subsection{\Q3: What is the ceiling, and why?}
A manifest that knows exactly which tools the user task needs (``oracle tools'') blocks \N{otoolBlk}\% of pairs
(\N{otoolBlkCI}\%). No tool-level policy, whether built by keywords, by an LLM or by hand per task, can do better on
these pairs, because the remaining \N{nUnblk} pairs (\N{unblockable}\%) are reachable with tools the user's task
legitimately needs. \N{nUnblkText} of them come from one travel injection that only asks the agent to recommend a hotel
in its answer, and needs no tool at all. Most of the rest are in banking (\N{unblockablebanking}\% of banking pairs).
There, seven of the nine injections need no sensitive tool other than \tool{send\_money}, and paying, refunding or transferring
is what many banking user tasks are about. Telling the user's payment from the attacker's would need checks on the
arguments (which IBAN, how much) or on where they came from, as in Progent and CaMeL
\cite{shi2025progent,debenedetti2025camel}.

The hand-written manifest reaches \N{manBlk}\% of this \N{otoolBlk}\% ceiling. The gap has two sources. In \N{gapCat}
pairs the user task needs one tool of a category and the attack uses another tool of the same category: paying a bill
opens \tool{money\_out}, which also contains \tool{update\_scheduled\_transaction}. Using oracle categories instead of
oracle tools lowers banking from \N{otoolBlkbanking}\% to \N{ocatBlkbanking}\%. In \N{gapKw} pairs a keyword in the
instruction opens a category the task does not need. In the other direction, \N{extraBlk} pairs are blocked by our
manifest but not by the oracle, because it also denies a tool the user task needs. Those blocks cost utility, so we
also report pairs that are both blocked and kept: \N{manBoth}\% (\N{manBothCI}\%), against \N{otoolBoth}\% for the
oracle.

\subsection{\Q4: End-to-end with a model}
\textcolor{red}{\textbf{[Pending. This subsection is filled only from the completed runs of
\texttt{code/run\_llm.py} (Claude Haiku; no defense, capscope, LLM tool filter; AgentDojo's
\texttt{important\_instructions} attack) via \texttt{code/analyze\_llm.py}. No number is entered by hand.]}}
